\section{Equipamiento y procedimiento}
\label{sec:equipamiento}

\subsection{Equipamiento utilizado}
En la \autoref{tab:equipamiento} se listan los instrumentos y elementos utilizados en las mediciones.

\begin{table}[H]
    \centering
    \caption{Equipamiento utilizado en las mediciones.}
    \label{tab:equipamiento}
    \begin{tabular}{p{4.2cm} p{4.6cm} p{5.6cm}}
        \toprule
        \textbf{Elemento} & \textbf{Marca y modelo} & \textbf{Características relevantes} \\
        \midrule
        Analizador de espectro & \pendiente{completar} & Rango de frecuencias, impedancia de entrada ($50 \, [\Omega]$), nivel máximo de entrada \\
        Bloqueador de continua (DC block) & \pendiente{completar} & Protege la entrada del analizador de tensiones continuas, por ejemplo la alimentación de un preamplificador de antena \\
        Antena para AM & \pendiente{completar} & \pendiente{tipo: hilo largo, cuadro o varilla} \\
        Antena para FM y TV & \pendiente{completar} & \pendiente{tipo: dipolo, Yagi o log-periódica; impedancia} \\
        Cables y conectores & \pendiente{completar} & \pendiente{tipo de cable (RG-58, RG-6), conectores BNC, N o F} \\
        Adaptador de impedancia & \pendiente{si se usó} & Adaptador de $75$ a $50 \, [\Omega]$ \\
        Televisor o receptor ISDB-Tb & \pendiente{completar} & Búsqueda automática de canales \\
        Osciloscopio & \pendiente{si se usó} & \\
        \bottomrule
    \end{tabular}
\end{table}

\subsection{Procedimiento}
\begin{enumerate}
    \item Se identificaron en el CABFRA \cite{cabfra} las bandas, los canales y las frecuencias de cada servicio, resumidos en la \autoref{tab:canalizacion}.
    \item Se conectó la antena correspondiente al analizador de espectro a través del bloqueador de continua, como se muestra en la \autoref{fig:conexion}. Solo se midieron señales recibidas: en ningún caso se conectó la salida de un transmisor al analizador.
    \item Antes de cada medición se configuró el instrumento y se registraron la frecuencia central o las de inicio y fin, el \textit{span}, el nivel de referencia, el ancho de banda de resolución (RBW), el ancho de banda de video (VBW), la atenuación, la preamplificación, la unidad y la impedancia de entrada.
    \item Se partió de los ajustes recomendados por la cátedra (nivel de referencia de $-30 \, [dBm]$, ganancia de $5$ a $10 \, [dB]$, RBW de $9 \, [kHz]$ y VBW de $100 \, [kHz]$) y se adaptaron a cada señal. Para explorar una banda completa conviene un RBW amplio, y para resolver la estructura de una emisión hace falta un RBW menor que los detalles que se quieren ver: las bandas laterales de AM, por ejemplo, requieren un RBW de $1 \, [kHz]$ o menos \cite{keysight_sa}.
    \item Se tomaron capturas de pantalla de cada medición, en las que se leen la señal y los parámetros, y fotografías de la antena, de los instrumentos y de las conexiones.
\end{enumerate}

\begin{table}[H]
    \centering
    \caption{Canalización de los servicios medidos según el CABFRA \cite{cabfra}.}
    \label{tab:canalizacion}
    \begin{tabular}{p{3.6cm} p{3.4cm} p{3.0cm} p{4.4cm}}
        \toprule
        \textbf{Servicio} & \textbf{Banda} & \textbf{Canalización} & \textbf{Observaciones} \\
        \midrule
        Radiodifusión sonora AM & $535$ a $1.605 \, [kHz]$ & 107 canales de $10 \, [kHz]$ & Se extiende con 10 canales más hasta $1.705 \, [kHz]$ (plan de Río de Janeiro, 1988) \\
        Radiodifusión sonora FM & $88$ a $108 \, [MHz]$ & Canales 201 a 300, de $200 \, [kHz]$ & El canal 201 está centrado en $88{,}1 \, [MHz]$ \\
        Televisión VHF baja & $54$ a $72$ y $76$ a $88 \, [MHz]$ & Canales 2 a 6, de $6 \, [MHz]$ & \\
        Televisión VHF alta & $174$ a $216 \, [MHz]$ & Canales 7 a 13, de $6 \, [MHz]$ & \\
        Televisión UHF & $470$ a $608$ y $614$ a $698 \, [MHz]$ & Canales 14 a 36 y 38 a 51, de $6 \, [MHz]$ & El canal 37 queda reservado a radioastronomía; la TDA usa esta banda \\
        \bottomrule
    \end{tabular}
\end{table}

El canal $N$ de FM está centrado en
\begin{equation}
    f_{FM}(N) = 88{,}1 \, [MHz] + 0{,}2 \, [MHz] \cdot (N - 201)
    \label{eq:canal_fm}
\end{equation}
y el canal UHF $N$ de televisión ocupa de $470 + 6\,(N-14)$ a $476 + 6\,(N-14) \, [MHz]$. En ISDB-Tb la portadora central se desplaza $1/7 \, [MHz]$ por encima del centro del canal, de modo que
\begin{equation}
    f_{ISDB}(N) = 473 + 6\,(N - 14) + \frac{1}{7} \, [MHz]
    \label{eq:canal_isdb}
\end{equation}

\begin{figure}[H]
    \centering
    \begin{tikzpicture}[node distance=1.6cm, every node/.style={font=\small},
        caja/.style={draw, rounded corners, minimum height=1cm, minimum width=2.4cm, align=center, fill=blue!6}]
        \node[caja] (ant) {Antena};
        \node[caja, right=of ant] (adap) {Adaptador\\$75/50 \, \Omega$\\(si corresponde)};
        \node[caja, right=of adap] (dc) {Bloqueador\\de continua};
        \node[caja, right=of dc, fill=orange!12] (sa) {Analizador\\de espectro\\$50 \, \Omega$};
        \draw[-{Stealth}] (ant) -- node[above]{coaxial} (adap);
        \draw[-{Stealth}] (adap) -- (dc);
        \draw[-{Stealth}] (dc) -- (sa);
    \end{tikzpicture}
    \caption{Conexión de la antena al analizador de espectro.}
    \label{fig:conexion}
\end{figure}

\begin{figure}[H]
    \centering
    \pendiente{Fotografía de la antena, de los instrumentos y de las conexiones (\texttt{Figuras/mediciones/montaje.jpg})}
    \caption{Montaje utilizado en el laboratorio.}
    \label{fig:montaje}
\end{figure}

\subsection{Impedancia y desadaptación}
El analizador de espectro tiene una entrada de $50 \, [\Omega]$, mientras que las antenas de televisión, los cables RG-6 y los receptores de radiodifusión trabajan normalmente en $75 \, [\Omega]$. Conectarlos directamente produce un coeficiente de reflexión
\begin{equation}
    \Gamma = \frac{Z_L - Z_0}{Z_L + Z_0} = \frac{50 - 75}{50 + 75} = -0{,}2
\end{equation}
y una pérdida por desadaptación de
\begin{equation}
    L_{des} = -10 \log\left(1 - |\Gamma|^2\right) = -10 \log(0{,}96) = 0{,}18 \, [dB]
\end{equation}
que es despreciable frente a la incertidumbre del resto de la medición. Un adaptador de mínima pérdida hace la transición sin reflexiones, pero introduce una atenuación conocida de unos $5{,}7 \, [dB]$ que debe sumarse a la lectura. En las mediciones \pendiente{se usó / no se usó} adaptador.
